%% kapa: exam=1 round=28 year=2017 period=1 form=A subject=civil_law subject_ko=민법 seq=14 number=14
%% source: Q-Net 감정평가사 기출문제 (cst00309) artlSeq=5199603 file=extracted/1008/5199603/2017년 제28회 감정평가사 1차시험 문제지/1교시 A형.pdf
%% conversion: kordoc 4.24.51 -> tools/gen1.py (auto); review=proofread-sonnet; flags=none
\begin{question}{14}
복대리에 관한 설명으로 옳지 않은 것은? (다툼이 있으면 판례에 따름)
\begin{choices}
  \item 복대리인은 그 권한 내에서 본인을 대리한다.
  \item 임의대리인은 본인의 승낙이 있거나 부득이한 사유 있는 때가 아니면 복대리인을 선임하지 못한다.
  \item 법정대리인이 부득이한 사유로 복대리인을 선임한 경우, 그 선임감독에 관한 책임만이 있다.
  \item 복대리인을 선임하더라도 대리인의 대리권은 소멸하지 않는다.
  \item 복대리인이 선임한 대리인은 모두 법정대리인이다.
\end{choices}
\end{question}
\answer{5}
